\documentclass[../main]{subfiles}
\begin{document}
\ifSubfilesClassLoaded{\setcounter{page}{62}}{}
\section{Висновок}
\label{sec:11_zaklyuchenie}

Лише тепер, коли ми в загальних рисах виклали основні моменти комунізму та показали ті паростки, з яких він може прорости як суспільний лад, ми вважаємо, що можемо перейти до визначень. В іншому випадку, вони були б мертвими, позбавленими сенсу абстракціями, а не відповіддю на питання, винесене в назву: «що таке комунізм?».

Комунізм (комуністичне суспільство) – це вільне об’єднання людей, які свідомо творять своє життя, свідомо покладаються на власну вільну діяльність.

Комунізм – це безпосередньо колективна форма вільної діяльності, як форма суспільства, що творить власний розвиток, яка стала тотальною.

Комунізм – це, як казали Маркс і Енгельс у «Маніфесті», «об’єднання, в якому вільний розвиток кожного є умовою вільного розвитку всіх».

Настав час знову повторити гасло, яке й досі не втратило своєї актуальності:

ПРОЛЕТАРІ ВСІХ КРАЇН, ОБ’ЄДНУЙТЕСЯ!
\end{document}
